%% =======================================================================
%  PREAMBULO — paquetes, estilo y macros
% =======================================================================

%% --- Idioma y codificacion ---------------------------------------------
\usepackage[utf8]{inputenc}
\usepackage[T1]{fontenc}
\usepackage[spanish,es-noquoting,es-tabla]{babel}
\usepackage{lmodern}
\usepackage{microtype}

%% --- Matematica --------------------------------------------------------
\usepackage{amsmath,amssymb,amsfonts}
\usepackage{bm}
\usepackage{siunitx}
\sisetup{
  output-decimal-marker = {,},
  per-mode              = symbol,
  range-units           = single,
  detect-weight         = true
}
\DeclareSIUnit\baud{Bd}
\DeclareSIUnit\dBm{dBm}
\DeclareSIUnit\dB{dB}
\DeclareSIUnit\ppm{ppm}

%% --- Figuras y tablas --------------------------------------------------
\usepackage{graphicx}
\graphicspath{{figuras/}}
\usepackage{booktabs}
\usepackage{multirow}
\usepackage{longtable}
\usepackage{float}
\usepackage[font=small,labelfont=bf,width=0.9\textwidth,
            justification=centering]{caption}
% Los epigrafes de las tablas van debajo de la tabla.
\captionsetup[table]{position=bottom}
\usepackage{subcaption}
\usepackage{tikz}
\usetikzlibrary{arrows.meta,positioning,calc,fit,backgrounds}


%% --- Pagina ------------------------------------------------------------
\usepackage[a4paper,inner=3cm,outer=2.5cm,top=2.5cm,bottom=2.5cm]{geometry}
% Parrafos sin sangria, separados por un espacio vertical.
\usepackage{parskip}
% Sin estirar el texto para llenar la pagina: el blanco sobrante queda
% al final y el interlineado no cambia de una pagina a otra.
\raggedbottom
\usepackage{fancyhdr}
\setlength{\headheight}{14.5pt}
\addtolength{\topmargin}{-2pt}
\pagestyle{fancy}
\fancyhf{}
\fancyhead[LE,RO]{\thepage}
\fancyhead[RE]{\nouppercase{\leftmark}}
\fancyhead[LO]{\nouppercase{\rightmark}}
\renewcommand{\headrulewidth}{0.4pt}

% Las paginas en blanco que \cleardoublepage mete para que cada capitulo
% abra en impar salian con el encabezado y el numero de pagina del capitulo
% anterior. Esto las deja realmente en blanco. Es la redefinicion estandar
% de book.cls con un \thispagestyle{empty} agregado.
\makeatletter
\renewcommand{\cleardoublepage}{%
  \clearpage
  \if@twoside
    \ifodd\c@page\else
      \null
      \thispagestyle{empty}%
      \newpage
      \if@twocolumn\null\newpage\fi
    \fi
  \fi}
\makeatother

%% --- Titulo de capitulo con el numero enmarcado -------------------------
%  Reproduce el estilo de la tesis de referencia. Si no te gusta,
%  comentar este bloque entero y queda el estilo estandar de book.
\usepackage{titlesec}
\titleformat{\chapter}[display]
  {\normalfont\Large\filleft}
  {\raisebox{0.6ex}{\normalsize\scshape \chaptertitlename}\hspace{0.4em}%
   \fbox{\parbox[c][2.1em][c]{2.1em}{\centering\Huge\thechapter}}}
  {0.6em}
  {\vspace{0.4em}\huge\filright}
\titlespacing*{\chapter}{0pt}{20pt}{40pt}

%  \paragraph viene en estilo corrido, con el texto pegado al titulo. Esto lo
%  pasa a bloque, con el titulo en su propio renglon. Tiene que ir despues del
%  \usepackage{titlesec} de arriba, porque de ahi sale \titleformat.
\titleformat{\paragraph}[block]
  {\normalfont\bfseries}{}{0pt}{}
\titlespacing*{\paragraph}{0pt}{1.5ex plus .5ex minus .2ex}{0.5ex}

%% --- Referencias cruzadas ----------------------------------------------
\usepackage{hyperref}
% Sin esto las URLs largas de la bibliografia entran como una caja
% indivisible. TeX las manda solas a su renglon y estira los espacios de las
% lineas anteriores para justificarlas. xurl deja cortarlas en cualquier
% punto y los huecos desaparecen.
\usepackage{xurl}
\hypersetup{
  colorlinks = true,
  linkcolor  = black,
  citecolor  = black,
  urlcolor   = blue!60!black,
  pdftitle   = {Receptor optico coherente DP-QPSK para enlaces LEO-GEO},
  pdfauthor  = {Ezequiel Brandan}
}
\usepackage[nameinlink,spanish]{cleveref}
% babel con es-tabla rotula los epigrafes como "Tabla", asi que cleveref
% tiene que decir lo mismo. Con el "cuadro" que trae por defecto salian
% referencias como "la cuadro 7.10", que ademas no concuerda en genero.
\crefname{table}{Tabla}{Tablas}
\Crefname{table}{Tabla}{Tablas}
\crefname{figure}{Figura}{Figuras}
\Crefname{figure}{Figura}{Figuras}
\crefname{equation}{Ecuaci\'on}{Ecuaciones}
\Crefname{equation}{Ecuaci\'on}{Ecuaciones}
\crefname{section}{Secci\'on}{Secciones}
\Crefname{section}{Secci\'on}{Secciones}
\crefname{subsection}{Secci\'on}{Secciones}
\Crefname{subsection}{Secci\'on}{Secciones}

% Vineta redonda en las listas (babel-spanish pone un guion).
\AtBeginDocument{\renewcommand{\labelitemi}{\textbullet}}

%% --- Utilidades de escritura -------------------------------------------
\usepackage{xcolor}
% Bloques comentados: \begin{comment} ... \end{comment}. Cada etiqueta tiene
% que ir sola en su linea y los bloques no se pueden anidar. El entorno lo da
% verbatim, que es parte de la base de LaTeX; el paquete comment hace lo mismo
% pero hay que instalarlo aparte.
\usepackage{verbatim}
% Antes de entregar: \newcommand{\todo}[1]{} y desaparecen todas.
\newcommand{\todo}[1]{\textcolor{red}{\textbf{[TODO: #1]}}}
\newcommand{\rev}[1]{\textcolor{blue}{#1}}

%% --- Lista de puntos sin sangria ----------------------------------------
% La vineta queda al mismo nivel que el margen del texto que la rodea, y el
% cuerpo del punto cuelga de ella. El \makelabel redefinido es lo que alinea
% la vineta a la izquierda de su caja: sin eso list la pega a la derecha y el
% punto arranca corrido. No hace falta enumitem.
\newenvironment{puntos}{%
  \begin{list}{$\bullet$}{%
    \setlength{\leftmargin}{1.5em}%
    \setlength{\labelwidth}{1em}%
    \setlength{\labelsep}{0.5em}%
    \setlength{\itemindent}{0pt}%
    \setlength{\listparindent}{0pt}%
    \setlength{\topsep}{\medskipamount}%
    \setlength{\parsep}{\parskip}%
    \setlength{\itemsep}{\medskipamount}%
    \renewcommand{\makelabel}[1]{##1\hfil}%
  }%
}{%
  \end{list}%
}

%% --- Entorno para la lista de acronimos ---------------------------------
\newenvironment{listaacronimos}
  {\begin{list}{}{%
     \setlength{\labelwidth}{2.4cm}
     \setlength{\leftmargin}{3.2cm}
     \setlength{\labelsep}{0.8cm}
     \setlength{\itemsep}{0.15em}
     \renewcommand{\makelabel}[1]{\textbf{##1}\hfil}}}
  {\end{list}}

%% =======================================================================
%  MACROS DEL TRABAJO
% =======================================================================
\newcommand{\Rs}{R_{s}}
\newcommand{\Ts}{T_{s}}
\newcommand{\Bref}{B_{\mathrm{ref}}}
\newcommand{\OSNR}{\mathrm{OSNR}}
\newcommand{\EVM}{\mathrm{EVM}}
\newcommand{\BER}{\mathrm{BER}}
\newcommand{\EbNo}{E_{b}/N_{0}}
\newcommand{\EsNo}{E_{s}/N_{0}}
\newcommand{\Pcw}{P_{\mathrm{cw}}}
\newcommand{\nsp}{n_{\mathrm{sp}}}
\newcommand{\Efield}{\mathbf{E}}
\newcommand{\Jones}[1]{\mathbf{#1}}
\newcommand{\fdot}{\dot{f}}
\newcommand{\dnu}{\Delta\nu}
\newcommand{\Df}{\Delta f}
\newcommand{\abs}[1]{\left\lvert #1 \right\rvert}
\newcommand{\E}[1]{\mathbb{E}\!\left[#1\right]}
% Etiquetas de las configuraciones de deterioros del capitulo de tolerancias
\newcommand{\Cfg}[1]{\mbox{C#1}}
\DeclareMathOperator{\sinc}{sinc}
\DeclareMathOperator{\Qfun}{Q}